\section{The actual arithmetic zero divisor on the retained inverse cover}
\label{sec:zt-zero-transport}

This section uses the arithmetic spectral coordinate $s$, the Newman time
$t$, and the original inverse-root coordinate $r$.  In particular, the
letter $s$ here is not the completed-zeta argument $\rho$.  The complete
coordinate dictionary is
\begin{equation}
 a=2s=-r^2,\qquad z=2s=-r^2,\qquad
 \rho=\tfrac12+is=\tfrac12-\tfrac i2r^2.
 \label{eq:zt-dictionary}
\end{equation}
The maps $s\mapsto a$ and $s\mapsto z$ have inverses $a\mapsto a/2$
and $z\mapsto z/2$; $s\mapsto\rho$ has inverse
$\rho\mapsto-i(\rho-1/2)$.  The last root map is a two-sheet map,
whose two local inverses on a simply connected open subset of
$\mathbb C_s\setminus\{0\}$ are $r=\sqrt{-2s}$ and
$r=-\sqrt{-2s}$.  No single inverse on all of $\mathbb C_s\setminus\{0\}$
is asserted.

The primary analytic conventions below are those of Brad Rodgers and
Terence Tao\cite[Section 1, equations (1)--(4)]{rodgersTao2020} and
D.~H.~J.~Polymath\cite[Section 3, Proposition 3.1]{polymath2019}.
The local calculations below prove their transport in the coordinates
\eqref{eq:zt-dictionary}, including the analytic factors in zero motion.

\subsection{The specified datum, its moments, and the branch point}

Keep the exact kernel and constants
\begin{align}
 \Phi(u)&=\sum_{n=1}^{\infty}
 (2\pi^2 n^4 e^{9u}-3\pi n^2e^{5u})e^{-\pi n^2e^{4u}},\notag\\
 H_t(z)&=\int_0^\infty e^{tu^2}\Phi(u)\cos(zu)\,du,\notag\\
 Z(t,s)&=8H_t(2s),\qquad
 f(t,r)=Z(t,-r^2/2)=8H_t(-r^2).
 \label{eq:zt-datum}
\end{align}
Thus $Z(0,s)=\Xi(s)=\xi(1/2+is)$ by the cited Fourier identity,
where $\xi(\rho)=\rho(\rho-1)\pi^{-\rho/2}
\Gamma(\rho/2)\zeta(\rho)/2$, with its entire continuation at the
removable exceptional points.  This identity fixes the initial datum;
the inverse cubic below is never substituted for it.

For $u\geq0$ each summand of $\Phi$ is positive, since its polynomial
factor is
\[
 \pi n^2 e^{5u}(2\pi n^2e^{4u}-3)>0.
\]
For a completely explicit integrable majorant put
$C_\Phi=\sum_{n\geq1}n^4e^{-\pi(n^2-1)}<\infty$.
The finiteness follows by comparison with
$\sum_{n\geq1}n^4e^{-\pi(n-1)}$.  With
$\lambda=\pi e^{4u}\geq\pi$, one has
\[
 0<\Phi(u)\leq2\pi^2 e^{9u}
 \sum_{n\geq1}n^4e^{-\lambda n^2}
 \leq2\pi^2C_\Phi e^{9u-\pi e^{4u}}.
\]
On $|t|\leq T$, $|s|\leq S$, every derivative of the integrand in
\eqref{eq:zt-datum} is bounded by a constant times
$u^N\exp(Tu^2+(9+2S)u-\pi e^{4u})$, for an integer $N\geq0$.
This is integrable: $\pi e^{4u}/2$ eventually exceeds
$Tu^2+(9+2S)u+N\log(1+u)$, and
$\exp(-\pi e^{4u}/2)$ is integrable.  The same argument applies on
compact $r$ sets, with $2S$ replaced by a bound for $|r|^2$.
Differentiation under the integral and locally uniform holomorphic
integration therefore show that $Z$ is entire on $\mathbb C_t\times
\mathbb C_s$, and $f$ is entire on $\mathbb C_t\times\mathbb C_r$.
They give the exact equations
\begin{equation}
 Z_t=-\tfrac14 Z_{ss},\qquad
 f_t=\mathcal L f,\qquad
 \mathcal L=-\frac1{4r^2}\partial_r^2+
                 \frac1{4r^3}\partial_r\quad(r\ne0).
 \label{eq:zt-heat}
\end{equation}
Indeed $f_r=-rZ_s$ and $f_{rr}=-Z_s+r^2Z_{ss}$, so the two terms
in $\mathcal Lf$ are $Z_s/(4r^2)-Z_{ss}/4$ and
$-Z_s/(4r^2)$, respectively.  Their sum is $-Z_{ss}/4$.

Define the actual moments, with their original measure,
\[
 M_k(t)=\int_0^\infty u^{2k}e^{tu^2}\Phi(u)\,du
 \quad(k=0,1,2,\ldots).
\]
The majorant just proved gives $M_k'(t)=M_{k+1}(t)$ and
\begin{align}
 \partial_s^{2k}Z(t,0)&=8(-1)^k2^{2k}M_k(t),&
 \partial_s^{2k+1}Z(t,0)&=0,\notag\\
 f(t,r)&=8\sum_{k=0}^{\infty}
          \frac{(-1)^kM_k(t)}{(2k)!}r^{4k}.
 \label{eq:zt-moments}
\end{align}
The series converges locally uniformly in both variables, by expansion
of the cosine and the same compact-set majorant.  In particular
\begin{align*}
 \Xi(0)&=8M_0(0)
     =-\frac{\Gamma(1/4)\zeta(1/2)}{8\pi^{1/4}}>0,\\
 \Xi'(0)&=0,\qquad \Xi''(0)=-32M_1(0)<0,\\
 f(t,r)&=8M_0(t)-4M_1(t)r^4+\tfrac13M_2(t)r^8+O(r^{12}).
\end{align*}
The equality with $\Gamma(1/4)\zeta(1/2)$ follows by evaluating the
specified completed-zeta identity, while both strict signs follow
directly from the positive integrals.  Neither sign requires a numerical
evaluation of $\zeta(1/2)$.

For every real $t$ and $v$ there is the stronger exact assertion, also
recorded in the first paragraph of Polymath's
Section 3\cite{polymath2019},
\begin{equation}
 Z(t,iv)=8\int_0^\infty e^{tu^2}\Phi(u)\cosh(2vu)\,du>0.
 \label{eq:zt-imaginary-axis-positive}
\end{equation}
In particular the actual arithmetic zero divisor avoids $s=0$ for every
real Newman time.  The completed root cover has a branch point there,
but no arithmetic zero reaches that branch point at any real time.
Evenness of the cosine also proves, for all complex $t,s,r$,
\begin{equation}
 Z(t,-s)=Z(t,s),\qquad f(t,-r)=f(t,r),\qquad f(t,ir)=f(t,r).
 \label{eq:zt-symmetries}
\end{equation}
For real $t$ complex conjugation gives
$\overline{Z(t,s)}=Z(t,\overline s)$ and
$\overline{f(t,r)}=f(t,\overline r)$, because the kernel is real.

\subsection{Exact function maps and the removable and nonremovable terms}

Write $\mathcal O(\mathbb C)$ for entire functions with the topology of
uniform convergence on compact sets.  The pullback
\[
 \pi^*: \mathcal O(\mathbb C_s)\longrightarrow\mathcal O(\mathbb C_r),
 \qquad q\longmapsto q(-r^2/2)
\]
is injective, since every $s\in\mathbb C$ has a root of $r^2=-2s$.
Its image is exactly the even entire functions.  To prove the latter
statement without a choice of square root, expand an even entire
function as $g(r)=\sum_{n\geq0}a_{2n}r^{2n}$ and define
$q(s)=\sum_{n\geq0}(-2)^na_{2n}s^n$.  This series converges for every
$s$: choose $R>\sqrt{2|s|}$ and apply the Cauchy bound
$|a_{2n}|\leq\max_{|r|=R}|g(r)|R^{-2n}$.  It gives
$q(-r^2/2)=g(r)$, and proves uniqueness by comparing coefficients.
It also proves continuity of the inverse on compact sets.  Every entire
cover function has the unique decomposition
\[
 g(r)=q_0(-r^2/2)+r q_1(-r^2/2).
\]
Here $q_0$ is recovered from $(g(r)+g(-r))/2$ and $q_1$ from
$(g(r)-g(-r))/(2r)$, with the latter expression extended at $r=0$
by its power series.  The sheet involution $\iota g(r)=g(-r)$ has
trace map $g\mapsto g+\iota g$ with kernel exactly the odd entire
functions.  These are maps of actual holomorphic functions, before any
arithmetic quotient is formed.

For every integer $n\geq0$ the operator in \eqref{eq:zt-heat} gives
\begin{equation}
 \mathcal L(r^n)=\frac{n(2-n)}4r^{n-4}.
 \label{eq:zt-monomials}
\end{equation}
Consequently, for $g(r)=\sum_{n\geq0}a_nr^n$, its Laurent expansion is
\[
 \mathcal Lg=\frac{a_1}{4r^3}-\frac{3a_3}{4r}
       +\sum_{n\geq4}\frac{n(2-n)}4a_nr^{n-4}.
\]
The terms $n=0,2$ are zero.  Thus $\mathcal Lg$ extends to an entire
function exactly when $a_1=a_3=0$.  Its value at zero in that case is
$-2a_4=-g^{(4)}(0)/12$.  In particular
$\mathcal L f(t,0)=8M_1(t)=f_t(t,0)$: the singular coordinate
coefficients cancel on the actual arithmetic datum.

The largest subspace of $\mathcal O(\mathbb C_r)$ on which every
iterate $\mathcal L^j$ remains entire is the even subspace.  For even
monomials successive applications of \eqref{eq:zt-monomials} reach
$r^0$ or $r^2$, which are annihilated, without producing a pole.
Equivalently, $\mathcal L^j\pi^*q=(-1/4)^j\pi^*q^{(2j)}$.
Conversely, let $a_n$ be the first nonzero odd coefficient of $g$.
When $n=4j+1$, applying $\mathcal L^{j+1}$ produces a nonzero multiple
of $r^{-3}$ from this term; when $n=4j+3$ it produces a nonzero
multiple of $r^{-1}$.  Every factor
$(n-4\ell)(2-n+4\ell)/4$ is nonzero for these odd exponents.
Distinct initial powers retain distinct exponents after the same
number of applications, so no cancellation with a higher power is
possible.  Even powers cannot contribute a pole.  This proves the
assertion and maximality.  It asserts stability of all differential
iterates, not existence of a heat evolution for every even entire datum.

There is an equivalent spacetime statement.  A function $g$ holomorphic
on a product of a time disk with a full disk about $r=0$, satisfying
$g_t=\mathcal Lg$ for $r\ne0$, is even in $r$.  Indeed successive
time differentiation gives $\mathcal L^jg=\partial_t^jg$ on the
punctured disk, and the right side is holomorphic at zero for every $j$.
The preceding first-odd-coefficient argument, which is local in $r$,
applies at every time.  The equation therefore descends there through
the proved even-function inverse to $q_t=-q_{ss}/4$.
On a connected punctured plane the kernel of $\mathcal L$ consists
exactly of $A+Br^2$: multiplying the equation by $4r^3$ gives
$rg_{rr}=g_r$, hence $(g_r/r)'=0$ and then $g=A+Br^2$.

\subsection{The original cubic, all finite fibres, and the zero divisor}

Keep the original map in spatial coordinates $(x,y,w)$:
\begin{align*}
 F_1&=(1+xy)^3w+y^2(1+xy)(4+3xy),\\
 F_2&=y+3x(1+xy)^2w+3xy^2(4+3xy),\\
 F_3&=2x-3x^2y-x^3w.
\end{align*}
Its retained polynomial is
\[
 P_{a,b,c}(r)=cr^3-2r^2+br-2a,\qquad
 \alpha=\tfrac12P'_{a,b,c}(r).
\]
For $x\ne0$ the inverse chart is
\begin{equation}
 x=\alpha^{-1},\qquad y=r-\alpha,\qquad
 w=5\alpha^2-3r\alpha-c\alpha^3,
 \qquad P_{a,b,c}(r)=0,\quad P'_{a,b,c}(r)\ne0.
 \label{eq:zt-original-inverse}
\end{equation}
Here is the needed exact verification.  Set $r=y+1/x$ and
$\alpha=1/x$.  Solving $F_3=c$ gives the displayed expression for
$w$.  Substitution in the other components gives
$F_2=4r+2\alpha-3cr^2$ and $F_1=r^2+r\alpha-cr^3$.
The equality $F_2=b$ is $2\alpha=P'(r)$, and after its substitution
$2(F_1-a)=P(r)$.  These identities prove both directions and recover
$r$ from the inverse image.  In particular a multiple root $P'=0$
has no finite point in this chart.  On $x=0$, direct substitution
gives $F(0,y,w)=(w+4y^2,y,0)$; there is the additional point
$(0,b,a-4b^2)$ precisely when $c=0$.

The arithmetic slice $a=2s$, $b=c=0$ now gives, with no omitted
coefficient,
\[
 P_{2s,0,0}(r)=-2r^2-4s,\qquad \alpha=-2r.
\]
Its root relation is exactly $s=-r^2/2$.  For $r\ne0$ the spatial
inverse is the explicit map
\begin{equation}
 \Gamma(r)=\left(-\frac1{2r},\;3r,\;26r^2\right),
 \qquad F(\Gamma(r))=(-r^2,0,0).
 \label{eq:zt-spatial-lift}
\end{equation}
The inverse root is $y+1/x=3r-2r=r$.  For $s\ne0$, its two
roots give two distinct finite points $\Gamma(r)$ and $\Gamma(-r)$,
and the complete original fibre also has the point $(0,0,2s)$.
For $s=0$ the polynomial has its double root at zero, both inverse
chart points are excluded by $\alpha=0$, and the sole finite fibre
point is $(0,0,0)$.  The completed root point $r=0$ is not silently
assigned the undefined spatial coordinate $x=-1/(2r)$.

For real $t$, every zero $s_*$ of the specified $Z(t,\cdot)$ is
nonzero by \eqref{eq:zt-imaginary-axis-positive}.  It therefore has
exactly two zero lifts $r_*$ and $-r_*$ on the retained cover, and
these have the finite spatial representatives in
\eqref{eq:zt-spatial-lift}.  The actual scalar
$Z(t,F_1(x,y,w)/2)$ vanishes on all three points of the original fibre
above $(2s_*,0,0)$.  Thus the odd sheet observable is retained by the
cover, and the extra $x=0$ fibre point is separately recorded.

Multiplicity is also transported exactly.  At any fixed complex time,
write $Z(t,s)=(s-s_*)^mA(s)$ at a zero, where $A(s_*)\ne0$.
For $r_*\ne0$ with $r_*^2=-2s_*$,
\[
 f(t,r)=\left[-\tfrac12(r-r_*)(r+r_*)\right]^m A(-r^2/2).
\]
The factor $[-(r+r_*)/2]^mA(-r^2/2)$ is a unit near $r_*$,
so the multiplicity is $m$ there and also at $-r_*$.  The local
coordinate $s-s_*=-(r-r_*)(r+r_*)/2$ has derivative $-r_*\ne0$;
its holomorphic inverse is the corresponding branch of
$r=\sqrt{-2s}$.  It gives isomorphisms of the local function rings
and of their quotients by the respective zero ideals.  At $s_*=0$
the exact ideal is instead
$(s^m)\mathcal O_{\mathbb C_r,0}=(r^{2m})$; the factor
$(-1/2)^mA(-r^2/2)$ is a unit.  Thus the multiplicity doubles there.
The quotient at that point is the rank-two module
\[
 \mathbb C\{r\}/(r^{2m})
   =\mathbb C\{s\}/(s^m)\ \oplus\ r\mathbb C\{s\}/(s^m),
 \qquad s=-r^2/2.
\]
Separating even and odd coefficients proves existence and uniqueness
of this decomposition; the trace has precisely its odd summand as
kernel.  These local divisor rings do not assert that differentiation
acts on a frozen global arithmetic spectral quotient.

For $r=u+iv$ with $u,v\in\mathbb R$, the original arithmetic point is
\begin{equation}
 s=\frac{v^2-u^2}{2}-iuv,\qquad
 \rho=\frac12+uv+i\frac{v^2-u^2}{2}.
 \label{eq:zt-critical-coordinate}
\end{equation}
It lies on $\operatorname{Re}\rho=1/2$ exactly when $uv=0$,
which is exactly the union of the real and imaginary $r$ axes.
For real $t$ there are no cover zeros on $u=v$ or $u=-v$, since
these lines map into the imaginary $s$ axis of
\eqref{eq:zt-imaginary-axis-positive}.  In particular an imaginary
root coordinate alone does not give an off-critical arithmetic zero;
the exact arithmetic test in this coordinate is $uv\ne0$ at a zero
of $f(0,r)$.

\subsection{Simple zero motion with every analytic-unit term}

Consider a point of the actual zero locus with $Z_s(t_0,s_0)\ne0$.
The holomorphic implicit function theorem applies because $Z$ is
jointly entire and this derivative is nonzero.  It gives a unique
holomorphic zero $s=\sigma(t)$ in a product of small disks.
Differentiating $Z(t,\sigma(t))=0$ and using
\eqref{eq:zt-heat} gives the exact velocity
\begin{equation}
 \sigma'(t)=\frac{Z_{ss}(t,\sigma(t))}{4Z_s(t,\sigma(t))}.
 \label{eq:zt-simple-base}
\end{equation}
More generally write a local factorization $Z=AQ$, where $A$ is a
nonvanishing holomorphic function of $(t,s)$ and $Q$ is the full
local zero factor under consideration.  The product rule gives
\begin{equation}
 Q_t=-\tfrac14Q_{ss}-\frac{A_s}{2A}Q_s
             -\left(\frac{A_t}{A}+\frac{A_{ss}}{4A}\right)Q.
 \label{eq:zt-unit-pde}
\end{equation}
All functions in this equality are evaluated on that product of
disks; dividing by $A$ is valid there.  At a simple root it gives
\begin{equation}
 \sigma'=\frac14\frac{Q_{ss}}{Q_s}
                      +\frac12\frac{A_s}{A}.
 \label{eq:zt-unit-velocity}
\end{equation}
If $Q=\prod_{j=1}^m(s-\sigma_j(t))$ has distinct roots at the time
in question, factor off $s-\sigma_j$ and differentiate the remaining
product.  This proves
\[
 \frac{Q_{ss}(\sigma_j)}{Q_s(\sigma_j)}
       =2\sum_{k\ne j}\frac1{\sigma_j-\sigma_k},\qquad
 \sigma_j'=\frac12\sum_{k\ne j}\frac1{\sigma_j-\sigma_k}
                    +\frac12\frac{A_s(t,\sigma_j)}{A(t,\sigma_j)}.
\]
For a single local root $Q=s-\sigma$, its entire velocity is
$A_s/(2A)$.  Removing the unit would remove this actual velocity.

Away from $r=0$, a lifted simple root obeys $r^2=-2\sigma$,
so differentiation and the computed spatial derivatives give
\begin{equation}
 r'=-\frac{\sigma'}r
    =-\frac1{4r}\frac{Z_{ss}}{Z_s}
    =\frac1{4r^2}\frac{f_{rr}}{f_r}-\frac1{4r^3}.
 \label{eq:zt-simple-cover}
\end{equation}
All ratios are evaluated at the moving zero.  For a factorization
$f=BV$ on a cover disk, the complete product equation is
\begin{align}
 V_t={}&-\frac1{4r^2}V_{rr}
       +\left(\frac1{4r^3}-\frac{B_r}{2r^2B}\right)V_r\notag\\
 &+\left(-\frac{B_{rr}}{4r^2B}
              +\frac{B_r}{4r^3B}-\frac{B_t}{B}\right)V,
 \label{eq:zt-cover-unit-pde}
\end{align}
and at a simple root
\[
 r'=\frac1{4r^2}\frac{V_{rr}}{V_r}
                    +\frac{B_r}{2r^2B}-\frac1{4r^3}.
\]
For the two complete sheet lifts of a single base zero, take
$V=r^2+2\sigma(t)$ and
$B(t,r)=-A(t,-r^2/2)/2$, where $Z=A(s-\sigma)$.
At a root $V_{rr}/V_r=1/r$.  Its term $1/(4r^3)$ cancels the
displayed $-1/(4r^3)$ exactly, while
$B_r/B=-rA_s/A$.  The result is
$r'=-A_s/(2rA)=-\sigma'/r$.
This computation retains both sheet roots and the analytic unit;
it identifies the cancellation rather than deleting either term.
The original spatial velocity is also explicit:
\[
 \frac{d}{dt}\Gamma(r(t))
 =\left(\frac{r'}{2r^2},\ 3r',\ 52rr'\right),\qquad
 \frac{d}{dt}F(\Gamma(r(t)))=(2\sigma',0,0).
\]

\subsection{A double zero: exact pair equations and the unit drift}

At an actual double zero $(t_0,s_0)$, put $h=t-t_0$ and write
\[
 Z(t_0,s_0+\delta)=a_2\delta^2+a_3\delta^3+O(\delta^4),
 \qquad a_2\ne0.
\]
The local zero factor can be expressed as
$Q=(s-M)^2-D^2$, where $M$ and $D^2$ are holomorphic in $t$,
$M(t_0)=s_0$, and $D^2(t_0)=0$.  A construction not depending on
the naming of the two roots is to take a small circle containing
only the double zero and no boundary zero.  The integrals
\[
 p_k(t)=\frac1{2\pi i}\int_{\text{circle}}
         s^k\frac{Z_s(t,s)}{Z(t,s)}\,ds\quad(k=1,2)
\]
are holomorphic and equal the first two root power sums with
multiplicity, by the residue computation for a local zero factor.
Thus $M=p_1/2$ and $D^2=p_2/2-M^2$ are holomorphic.  The two
zeros, counted with multiplicity, are $M+D$ and $M-D$.
The holomorphic unit can itself be constructed on a smaller concentric
disk by the contour formula
\[
 A(t,s)=\frac1{2\pi i}\int_{\text{circle}}
       \frac{Z(t,\zeta)}{Q(t,\zeta)(\zeta-s)}\,d\zeta.
\]
The denominator is nonzero on the contour for all the retained times,
so the integral is jointly holomorphic.  At each fixed time $Z/Q$
extends holomorphically across its two zeros, counted with their
multiplicities, by their one-variable Taylor factorizations.  Cauchy's
formula therefore identifies the displayed integral with $Z/Q$, proving
$Z=AQ$.  At $t=t_0$ it is $Z(t_0,s)/(s-s_0)^2$, nonzero at $s_0$,
so $A$ is nonvanishing after shrinking the disks.

Put $b(t,s)=A_s(t,s)/A(t,s)$.  At every nearby time with $D\ne0$
the proved simple-zero equation gives the exact pair dynamics
\begin{align}
 (M+D)'&=\frac1{4D}+\tfrac12 b(t,M+D),\notag\\
 (M-D)'&=-\frac1{4D}+\tfrac12 b(t,M-D),\notag\\
 M'&=\tfrac14\bigl[b(t,M+D)+b(t,M-D)\bigr],\notag\\
 (D^2)'&=\tfrac12+\tfrac D2
              \bigl[b(t,M+D)-b(t,M-D)\bigr].
 \label{eq:zt-double-exact}
\end{align}
The sum and $D$ times the difference are even holomorphic functions
of $D$, hence holomorphic functions of $D^2$.  These equations extend
to $h=0$, and there give
\[
 (D^2)'(t_0)=\tfrac12,\qquad
 M'(t_0)=\frac{a_3}{2a_2}
        =\frac{Z_{sss}(t_0,s_0)}{6Z_{ss}(t_0,s_0)}.
\]
Indeed $A(t_0,s_0)=a_2$ and $A_s(t_0,s_0)=a_3$.
The full equation \eqref{eq:zt-unit-pde} still carries $A_t$ and
$A_{ss}$; their factors multiply $Q$ and vanish only when evaluated
at a zero.  They have not been assumed to vanish on the neighborhood.

The following explicit expansion proves the pair's first unit term
and its error order.  The heat equation gives the terms of weighted
degree at most three, with $\delta$ of weight one and $h$ of weight
two:
\[
 Z(t_0+h,s_0+\delta)
 =a_2(\delta^2-h/2)+a_3(\delta^3-3h\delta/2)
       +O((|\delta|+|h|^{1/2})^4).
\]
Set $h=\varepsilon^2$.  Dividing by $\varepsilon^2$ after
$\delta=\varepsilon y$ gives a holomorphic function of
$(\varepsilon,y)$ whose value at $\varepsilon=0$ is
$a_2(y^2-1/2)$.  Its two simple roots are
$y=\pm1/\sqrt2$, so the holomorphic implicit function theorem
gives two holomorphic $y$ branches.  Their next coefficient is found
by substituting $\delta=\lambda\varepsilon+\mu\varepsilon^2$:
$\lambda^2=1/2$ and
$2a_2\lambda\mu+a_3(\lambda^3-3\lambda/2)=0$.
Hence $\mu=a_3/(2a_2)$ and
\begin{equation}
 \sigma_\pm(t)=s_0\mathbin\pm\sqrt{h/2}
             +\frac{a_3}{2a_2}h+O(h^{3/2}).
 \label{eq:zt-double-expansion}
\end{equation}
This means convergent holomorphic expansions in $\varepsilon$,
with a chosen $\varepsilon^2=h$; reversing $\varepsilon$ exchanges
the root labels.  It follows that, on either punctured time branch,
\[
 \sigma_\pm'=\mathbin\pm\frac1{2\sqrt{2h}}
                 +\frac{a_3}{2a_2}+O(h^{1/2}).
\]
For real $t_0,s_0$ the coefficients are real.  For small positive
$h$ both branches are real, by uniqueness and conjugation of their
real $\varepsilon$ power series.  For small negative real $h$ their
imaginary parts are $\pm\sqrt{|h|/2}+O(|h|^{3/2})$, so they are a
nonreal conjugate pair.  This describes an actual double-zero event
wherever it occurs; it does not insert such an event at an unspecified
real arithmetic time.

At real time $s_0\ne0$ by \eqref{eq:zt-imaginary-axis-positive}.
Choose either $r_0^2=-2s_0$.  The inverse branch has
$dr/ds=-1/r_0$ and $d^2r/ds^2=-1/r_0^3$.  Taylor substitution
in \eqref{eq:zt-double-expansion} yields the two roots near that
particular sheet point,
\begin{equation}
 r_\pm(t)=r_0\mathbin\mp\frac1{r_0}\sqrt{h/2}
   -\left(\frac{a_3}{2a_2r_0}+\frac1{4r_0^3}\right)h
   +O(h^{3/2}).
 \label{eq:zt-double-root-expansion}
\end{equation}
The other sheet consists exactly of their negatives.  Both the
arithmetic unit drift and the curvature of the original root map
occur in its displayed coefficient of $h$.

\subsection{Every multiplicity and its interaction with ramification}

At an actual zero of finite multiplicity $m\geq1$ write
$Z(t_0,s_0+\delta)=\sum_{k\geq m}a_k\delta^k$, $a_m\ne0$.
There is finite multiplicity because $Z(t_0,\cdot)$ is not the zero
function.  For real $t_0$ positivity at zero already proves this.
For complex $t_0$ put
$K(u)=e^{t_0u^2}\Phi(|u|)$ on $\mathbb R$.  The proved decay shows
that $K$ is continuous, bounded, and integrable, and $K(0)=\Phi(0)>0$.
Its Fourier transform satisfies
$\widehat K(\xi)=2H_{t_0}(\xi)$ for real $\xi$, by evenness.
If $Z(t_0,\cdot)$ were identically zero, this transform would vanish.
For $\epsilon>0$, the Gaussian integral gives
\[
 G_\epsilon(x)=\frac{e^{-x^2/(4\epsilon)}}{\sqrt{4\pi\epsilon}}
   =\frac1{2\pi}\int_{\mathbb R}e^{-\epsilon\xi^2}e^{ix\xi}\,d\xi.
\]
Fubini is justified by $K\in L^1$ and the integrable Gaussian in
$\xi$; it gives
$(K*G_\epsilon)(0)=(2\pi)^{-1}\int\widehat K(\xi)
e^{-\epsilon\xi^2}\,d\xi=0$.
On the other hand $(K*G_\epsilon)(0)\to K(0)>0$ as
$\epsilon\downarrow0$: on $|u|<\delta$ continuity bounds
$|K(u)-K(0)|$ arbitrarily small, while on its complement boundedness
of $K$ and the vanishing Gaussian tail bound the remaining integral.
This contradiction proves the required nonidentity with the zero
function for every complex time.

Repeated differentiation of the heat equation gives
$\partial_t^\ell Z=(-1/4)^\ell\partial_s^{2\ell}Z$.
The bivariate Taylor coefficient of $h^\ell\delta^j$ is therefore
\[
 \frac{(-1/4)^\ell}{\ell!}\frac{(j+2\ell)!}{j!}
 a_{j+2\ell}.
\]
Set $h=\varepsilon^2$, $\delta=\varepsilon y/\sqrt2$, with the
inverse substitutions retained whenever $\varepsilon\ne0$.
The holomorphic Taylor series is divisible by $\varepsilon^m$;
grouping the terms of degree $m$ and $m+1$ gives
\begin{align}
 \frac{Z(t_0+\varepsilon^2,s_0+\varepsilon y/\sqrt2)}
      {a_m(\varepsilon/\sqrt2)^m}
 ={}&\operatorname{He}_m(y)
 +\frac{a_{m+1}}{a_m}\frac{\varepsilon}{\sqrt2}
              \operatorname{He}_{m+1}(y)+O(\varepsilon^2),\notag\\
 \operatorname{He}_m(y)={}&
 \sum_{\ell=0}^{\lfloor m/2\rfloor}
 \frac{m!(-1)^\ell}{2^\ell\ell!(m-2\ell)!}y^{m-2\ell}.
 \label{eq:zt-all-multiplicity}
\end{align}
The quotient extends holomorphically at $\varepsilon=0$; its
remainder is uniform on each bounded $y$ set.  This follows directly
from the convergent bivariate Taylor series after substitution, so
the error does not assume a finite polynomial initial datum.

For completeness the polynomial has $m$ distinct real roots.
Differentiating its displayed finite sum proves
$\operatorname{He}_m'=m\operatorname{He}_{m-1}$ and
$\operatorname{He}_{m+1}=y\operatorname{He}_m-m\operatorname{He}_{m-1}$.
Starting with $\operatorname{He}_0=1$ and
$\operatorname{He}_1=y$, induction shows that the roots of two
successive polynomials strictly interlace: at consecutive roots of
$\operatorname{He}_m$, the values
$-m\operatorname{He}_{m-1}$ alternate sign, providing one root of
$\operatorname{He}_{m+1}$ between each consecutive pair.  At the
largest root, $\operatorname{He}_{m-1}>0$, so
$\operatorname{He}_{m+1}<0$ there and is positive sufficiently far
to the right.  At the smallest root, the sign is
$(-1)^m$ and the far-left sign of its monic degree $m+1$
polynomial is $(-1)^{m+1}$, providing a further root to the left.
There are thus $m+1$ distinct roots.  The induction begins with the
single root of $\operatorname{He}_1$ and gives the claimed reality,
simplicity, and interlacing.  Let these roots be $\lambda_j$.
At a root of $\operatorname{He}_m$ the recurrences also give
\[
 \frac{\operatorname{He}_{m+1}(\lambda_j)}
      {\operatorname{He}_m'(\lambda_j)}=-1.
\]
Apply the holomorphic implicit function theorem to
\eqref{eq:zt-all-multiplicity} at each $(0,\lambda_j)$.  The resulting
roots are
\begin{equation}
 \sigma_j(t)=s_0+\lambda_j\sqrt{h/2}
                   +\frac{a_{m+1}}{2a_m}h+O(h^{3/2}).
 \label{eq:zt-general-roots}
\end{equation}
They are all the nearby roots: a small fixed circle about $s_0$
has no boundary zero at $t_0$, and continuity followed by Rouch\'e's
theorem keeps its zero count at $m$; the $m$ disjoint roots just
constructed account for that count.  Writing
$Z(t_0,s)=(s-s_0)^mA_0(s)$ gives
$a_{m+1}/(2a_m)=A_0'(s_0)/(2A_0(s_0))$.  Thus the first common
drift at every multiplicity is exactly the initial analytic-unit
derivative.  Higher coefficients remain determined by the full Taylor
series and the exact unit equation \eqref{eq:zt-unit-pde}.

For $s_0\ne0$, on either chosen sheet $r_0^2=-2s_0$ the same inverse
coordinate Taylor expansion proves
\[
 r_j(t)=r_0-\frac{\lambda_j}{r_0}\sqrt{h/2}
   -\left(\frac{a_{m+1}}{2a_mr_0}
                      +\frac{\lambda_j^2}{4r_0^3}\right)h
   +O(h^{3/2}).
\]
At $s_0=0$, an event of the actual family is possible only at
nonreal complex time, by \eqref{eq:zt-imaginary-axis-positive}.
On this locus the exact evenness $Z(t,-s)=Z(t,s)$ forces
$m=2k$ and $a_{m+1}=0$.  Its Hermite roots are nonzero, since
$\operatorname{He}_{2k}(0)=(-1)^k(2k)!/(2^kk!)$.
Set $h=\eta^4$ and $\varepsilon=\eta^2$.  Equation
\eqref{eq:zt-general-roots} now gives
\[
 \sigma_j(t_0+\eta^4)=\frac{\lambda_j}{\sqrt2}\eta^2+O(\eta^6),
 \qquad
 r_{j,\pm}(\eta)=\mathbin\pm\eta
                 \sqrt{-\sqrt2\lambda_j}+O(\eta^5).
\]
The square root of each nonzero displayed constant has two choices.
To justify the expansion, $-2\sigma_j(t_0+\eta^4)/\eta^2$
extends holomorphically and has nonzero value
$-\sqrt2\lambda_j$ at zero; either holomorphic square root of that
unit gives exactly $r_{j,\pm}=\pm\eta\sqrt{-2\sigma_j/\eta^2}$.
The $2m$ lifted roots account for the cover multiplicity $2m$ proved
above.  This is the precise fourth-root time ramification over an
arithmetic multiple zero at the cover branch locus.  It does not
occur there at real Newman time, and it has no finite spatial inverse
chart at its central point because $P'=0$.

\paragraph{Verification status.}
The accompanying exact rational-polynomial checker verifies the chain
rule coefficients, monomial pole coefficients, original spatial lift,
analytic-unit product identities, and finite coefficients of the
collision expansions.  The convergence, positivity, local holomorphic
factorizations, root counts, and complete multiplicity statements are
proved in the text; the checker is not a substitute for those proofs.
No off-critical zero of the distinguished arithmetic initial function
is asserted by these calculations.
